\documentclass[11pt,a4paper]{article}
\usepackage[utf8]{inputenc}
\DeclareUnicodeCharacter{2248}{\ensuremath{\approx}}
\usepackage[T1]{fontenc}
\usepackage[brazil]{babel}
\usepackage{lmodern}
\usepackage[margin=2cm]{geometry}
\usepackage{amsmath,amssymb,bm}
\usepackage{graphicx,booktabs,float,longtable,url}
\usepackage{xr-hyper}
\usepackage[hidelinks]{hyperref}
% as secoes da formulacao (F-) e do relatorio das etapas (E-): compile os dois antes (o atualiza.sh faz isso)
\externaldocument[F-]{../../../docs/confiabilidade_campos/formulacao/formulacao}
\externaldocument[E-]{../ETAPAS1a3/relatorio_etapas}
% a data e o commit vem do atualiza.sh; compilado a mao, fica so a data
\providecommand{\versao}{17/11/2026}
\title{Confiabilidade de taludes com campos aleatórios no NeoPZ\\\large Situação: plano de trabalho, resultados
principais, hipóteses e próximos passos}
\date{\versao}
\begin{document}
\maketitle
\tableofcontents

\section{Objetivo e onde estamos}
O trabalho serve a dois artigos:
\begin{itemize}
  \item \textbf{Amostragem por importância (AI)} para taludes com campos aleatórios no MEF. A validação é no Ceron (2025),
    no rebaixamento, e no Cho (2010), um talude seco com $c'$ e $\varphi'$ como campos correlacionados e $P_f$ de 0,4\,\% a
    2\,\%.
  \item \textbf{Chuva em talude não saturado} (Rahardjo et al., 2007), com as mesmas ferramentas: Monte Carlo, FORM e AI.
\end{itemize}
O plano de trabalho (\path{Projects/PlanoDeTrabalho/README.md}) valida as ferramentas uma de cada vez, antes de usá-las
na chuva. As etapas 0 a 3 estão rodadas, e a 2c (gradiente modal) foi testada no PC, no SRM e no GIM. A etapa 4 (FORM e
AI na chuva) tem a referência pronta, e a etapa 5 (a investigação da chuva) ainda não começou. Tudo roda no cluster
fermi (4 nós, 48 processos); o PC serve para testes curtos.

Três documentos:
\begin{itemize}
  \item a \textbf{formulação}, citada como ``formulação, Seção $n$'':\newline
    \path{docs/confiabilidade_campos/formulacao/formulacao.pdf};
  \item o \textbf{relatório das etapas}, com todos os resultados, as tabelas e as figuras, citado como ``etapas, Seção
    $n$'':\newline \path{Projects/PlanoDeTrabalho/ETAPAS1a3/relatorio_etapas.pdf};
  \item este, com a situação, o resumo dos resultados, as hipóteses e os próximos passos. A Seção~\ref{sec:mapa} lista
    todos os documentos.
\end{itemize}

\section{O plano de trabalho e a situação}
\begin{table}[H]\centering\small
\resizebox{\ifdim\width>\textwidth\textwidth\else\width\fi}{!}{%
\begin{tabular}{lp{0.5\textwidth}p{0.35\textwidth}}
\toprule
Etapa & O quê & Situação \\
\midrule
0 & Cho2010 no fermi: direto, AI pela banda e pelo FORM & pronta \\
1 & SRM corrigido ($F_0 = 0{,}8$): o antigo contra o novo nas mesmas realizações & pronta (etapas, Seção~\ref{E-sec:srm}) \\
2 & FORM ($M = 10$, 20) e AI pelo ponto de projeto; teste linear sem MEF & pronta; sem viés (etapas, Seção~\ref{E-sec:vies}) \\
2b & base de Nyström contra a analítica & o mesmo ponto de projeto (etapas, Seção~\ref{E-sec:pp}); fase 4 de Nyström parcial \\
2c & gradiente do FORM pela projeção modal & validado no SRM e no GIM (etapas, Seção~\ref{E-sec:validacao}); fase 4 parcial \\
3 & $P_f$ pequena ($r = -0{,}7$), tamanho do piloto, custo & o direto vai de 36\,385 (129 falhas) a 110\,000 (CoV 5\,\%), com os 4 nós do fermi \\
4 & FORM e AI na chuva & o direto de 10\,000 terminou ($P_f = 5{,}51\,\%$); falta o código do FORM na chuva \\
5 & investigação da chuva & não começou \\
\bottomrule
\end{tabular}}
\end{table}

\section{Resultados principais}\label{sec:resultados}
Cada item cita a seção do relatório das etapas com o resultado completo.
\begin{itemize}
  \item \textbf{SRM corrigido} (etapas, Seção~\ref{E-sec:srm}). No Cho2010, as falhas são as mesmas nas duas versões
    (22 e 22 em 1000); na chuva, 4 realizações passam a
    falhar e 1 deixa de falhar (McNemar $p = 0{,}38$). O efeito é pequeno, e o SRM novo, o correto, fica.
  \item \textbf{AI contra o direto, $r = -0{,}5$} (etapas, Seção~\ref{E-sec:aidireto}). O direto de 24\,000 dá $P_f = 1{,}70\,\%$ (Cho: 1,71\,\%). As AI completas ficam a menos de 1{,}4 desvio dele, com ganho de até 16$\times$; só a direção errada da malha Tri(0) fica mais longe.
  \item \textbf{$P_f$ pequena, $r = -0{,}7$} (etapas, Seção~\ref{E-sec:r07}). O direto de 36\,385 dá
    $P_f = 0{,}35\,\%$ (129 falhas; Cho: ≈0,4\,\%). As AI pelo FORM acertam,
    com ganho de até 51$\times$.
  \item \textbf{FORM e as duas bases} (etapas, Seções~\ref{E-sec:nystrom} e~\ref{E-sec:pp}). O $\beta$ converge para
    $\approx 2{,}00$ nas duas bases, e a de Nyström chega lá com menos termos. O FORM sozinho superestima a $P_f$ em
    ≈30\,\%: ele serve para o deslocamento da AI. O ponto de projeto é o mesmo nas duas bases (a 3{,}2$^\circ$).
  \item \textbf{O gradiente modal} (etapas, Seção~\ref{E-sec:validacao}). Uma análise por iteração, no lugar de $4M + 1$.
    No SRM, fica a 9 a 11$^\circ$ das diferenças centrais; no GIM, a 3{,}1 e 5{,}5$^\circ$. O FORM do GIM converge em 5 análises
    para $\beta = 2{,}131$ ($M = 10$) e $2{,}026$ ($M = 20$), contra 2{,}115 e 2{,}012 das diferenças centrais no SRM.
  \item \textbf{Sem viés} (etapas, Seção~\ref{E-sec:vies}). As falhas do direto ficam todas onde a AI amostra, a
    variância verdadeira é a estimada, e as sementes independentes acertam o direto. A fase 4 parcial (etapas,
    Seção~\ref{E-sec:fase4}) dá 1{,}95 $\pm$ 0{,}15\,\% (8 sementes, 529 realizações) na KL; 1{,}80 $\pm$ 0{,}13\,\% (9 sementes, 564 realizações) em Nyström, contra 1{,}70\,\% do direto.
  \item \textbf{Variável de controle} (etapas, Seção~\ref{E-sec:cv}). Com o substituto ajustado no direto inteiro, de
    1{,}7 a 3{,}2 vezes menos realizações; com o ajustado na própria
    amostra (o que se tem na prática), de 1{,}4 a 2{,}1. Sem viés,
    mas o desvio que ela estima saiu de 15 a 26\,\% menor que o
    verdadeiro, nos lotes do direto.
  \item \textbf{Piloto e custo} (etapas, Seção~\ref{E-sec:custo}). Um piloto de 250 basta, e ele pode rodar em malha
    grossa; o FORM, não. Para CoV de 10\,\%, a AI pelo ponto de projeto custa de 6 a 29 vezes menos análises que
    o direto, já contando o FORM.
  \item \textbf{A referência da chuva} (etapas, Seção~\ref{E-sec:chuva}). O direto de 10\,000 dá
    $P_f = 5{,}51\,\%$, e a AI pela banda concorda com ele (1{,}4 desvio).
\end{itemize}

\section{Hipóteses}\label{sec:hip}
\subsection*{Confirmadas}
\begin{itemize}
  \item A AI não tem viés com qualquer deslocamento fixado antes; o deslocamento só muda a variância (teste linear e Ceron).
  \item O SRM antigo tinha um buraco logo abaixo de FS = 1; o efeito na $P_f$ é pequeno.
  \item \textbf{Há um só modo de ruptura} no Cho2010, nas duas bases: todas as falhas do direto ficam onde a AI pelo
    FORM amostra ($t = \bm\xi\cdot\bm\xi^*/\beta \ge 0{,}80$ na direção modal, com só
    3 das 408 abaixo de 1,5;
    $t \ge 2{,}18$ com $r = -0{,}7$), e a variância verdadeira da AI,
    calculada pelas falhas do direto, é a que a AI estima (etapas, Seção~\ref{E-sec:diag}). A banda de ruptura das falhas
    é a mesma no direto e na AI.
  \item \textbf{O ponto de projeto é o mesmo nas duas bases:} levado para as coordenadas de Nyström, o ponto analítico
    fica a 3{,}2$^\circ$ do de Nyström, e os campos de $\varphi'$ diferem em
    menos de 0,1$^\circ$ (etapas, Figura~\ref{E-fig:pp}).
  \item \textbf{A variável de controle dá um ganho extra sem custo:} de 1{,}7 a
    3{,}2 nas AI de $r = -0{,}5$ com o substituto do direto, e de 1{,}4 a 2{,}1 com o da própria amostra. No direto, a referência vai
    a 1{,}705 $\pm$ 0{,}040\,\%.
  \item \textbf{O piloto pode ser pequeno:} com 250 realizações, o deslocamento da banda custa só ≈5\,\% de variância.
  \item \textbf{O piloto pode rodar em malha grossa:} o elite é escolhido por ordem, e a ordem das realizações quase não
    muda. Com o piloto na Tri(0) (300 vezes mais barata), a AI pela banda acertou o direto ($-0{,}2$ desvio), com o mesmo
    CoV do piloto na malha usual.
  \item \textbf{O gradiente modal dá a direção do FORM com 1 análise,} a menos de ≈9 a 11$^\circ$ das diferenças centrais
    no SRM e de 3 a 6$^\circ$ no GIM, e o FORM converge em 3 a 5 análises.
  \item \textbf{A base de Nyström precisa de menos termos} que a analítica para o mesmo $\beta$.
  \item \textbf{A AI pelo FORM não tem viés} ($r = -0{,}5$), nas duas bases (etapas, Seção~\ref{E-sec:vies}). Com a semente 3, as AI saem todas em ≈1,55\,\% (os sorteios são os mesmos). As 5 sementes independentes da KL dão, em média, 1{,}728\,\% contra 1{,}700\,\% do direto e, com a variável de controle, 1{,}679 $\pm$ 0{,}034\,\% contra 1{,}705 $\pm$ 0{,}040\,\% ($-0{,}5$ desvio). A direção errada da Tri(0) é a única que subestima o próprio CoV (9{,}3 contra 12{,}3\,\%); com o CoV verdadeiro, o seu desvio cai de $-2{,}2$ para $-1{,}5$.
  \item \textbf{A AI pelo FORM com $r = -0{,}7$ também não tem viés detectável:} as sementes 3, 4 e 5, corrigidas, dão 0{,}352, 0{,}332, 0{,}365\,\%, contra 0{,}372 $\pm$ 0{,}020\,\% do direto corrigido (36\,385 realizações, 129 falhas): $-0{,}9$ desvio.
\end{itemize}
\subsection*{Descartadas}
\begin{itemize}
  \item \textbf{O FORM na malha Tri(0)} (76 equações, FS 1,47 contra 1,20): $\beta = 4{,}08$ e direção a 31,7$^\circ$ da
    certa. O erro da malha desloca o ponto de projeto, e não só o FS. Com a direção da Tri(0) e a distância da malha usual
    ($\beta = 2{,}32$), a AI funciona, mas pior que o FORM na malha usual.
\end{itemize}
\subsection*{Sob investigação}
\begin{itemize}
  \item \textbf{O campo de $k$ na chuva:} o gradiente modal pega $c'$ e $\varphi'$; o $k$ age pelo transiente e
    precisaria do adjunto. Hoje o $k$ pesa pouco na chuva f50 ($\alpha^2_k \approx 0$), e a AI continua sem viés sem
    deslocá-lo.
  \item \textbf{O desvio estimado pela variável de controle} saiu otimista nos lotes do direto (etapas,
    Seção~\ref{E-sec:cvlocal}); as sementes da fase 4 vão dizer se é acaso.
\end{itemize}

\section{Próximos passos}
\begin{enumerate}
  \item \textbf{Direto de $r = -0{,}7$ até 110\,000} (CoV de 5\,\%; com os 4 nós do fermi). Está em 36\,385 realizações. A fase 4 fica pausada até ele acabar.
  \item \textbf{Fase 4 no fermi} (≈8\,h): 10 sementes do FORM modal analítico e 10 do de Nyström, as mesmas nas duas
    bases. Dão a distribuição do estimador e fecham a questão do viés (etapas, Seção~\ref{E-sec:vies}). As partes que
    rodaram antes do cancelamento já estão juntas (etapas, Seção~\ref{E-sec:fase4}).
  \item \textbf{Etapa 4:} o FORM com gradiente modal e a AI pelo ponto de projeto na chuva, contra o direto de
    10\,000.
  \item \textbf{Variável de controle realista} (etapas, Seção~\ref{E-sec:cvlocal}): o substituto ajustado na própria
    amostra ($K = 50$, 5 dobras) em todas as AI, inclusive pela banda e na chuva; conferir o desvio estimado com as 10
    sementes da fase 4. Para o artigo da AI: o ganho contra o custo total (FORM, piloto e controle), de
    $P_f \approx 2\,\%$ a $0{,}4\,\%$.
  \item \textbf{O FORM da chuva pelo GIM?} O gradiente modal do GIM é mais preciso que o do SRM (5{,}5 contra 11{,}2$^\circ$ no mesmo ponto), e a região de
    falha é a mesma. Testar se vale na chuva, com as pressões do transiente.
\end{enumerate}

\section{Mapa dos documentos}\label{sec:mapa}
É a mesma lista do índice do \path{Projects/PlanoDeTrabalho/README.md}: as duas saem de
\path{Projects/PlanoDeTrabalho/documentos.py}. Os caminhos relativos são a partir de \path{~/projects/neopz-master}.
\begin{longtable}{p{0.44\textwidth}rp{0.46\textwidth}}
\toprule
Documento & p. & O que tem \\
\midrule
\endhead
\multicolumn{3}{p{0.96\textwidth}}{\textbf{Formulação: o solo não saturado e o fator de segurança (docs/)}} \\
\path{docs/nao_saturado/biot_nao_saturado_agua_ar5.pdf} & 49 & formulação poroelástica não saturada (esqueleto, água e ar), a implementação no NeoPZ, Rahardjo et al. (2007) e benchmarks; a versão 4 está ao lado \\
\path{docs/nao_saturado/formulacao/formulacao.pdf} & 25 & formulação acoplada u--pw--pa com Mohr--Coulomb associado: discretização, resíduos e jacobiano \\
\path{docs/nao_saturado/resumo/resumo_nao_saturado.pdf} & 24 & resumo da formulação u--pw--pa \\
\path{docs/nao_saturado/relatorio_discussao.pdf} & 7 & contraexemplos a Rahardjo et al. (2007): discussão dos resultados, argumentos e verificação \\
\path{docs/nao_saturado/artigo/artigo.pdf} & 19 & artigo: a intensidade crítica de chuva em taludes não saturados, e quanto as análises desacopladas e com ar passivo a superestimam \\
\path{docs/talude_srm/slope_mohr_coulomb/relatorio.pdf} & 18 & talude Mohr--Coulomb no NeoPZ: correções da plasticidade, FS e rebaixamento \\
\path{docs/talude_srm/slope_mohr_coulomb/nota_uqp.pdf} & 6 & nota: formulação de três campos u--q--p \\
\path{docs/talude_srm/fs_arc_length/relatorio_fs.pdf} & 34 & FS de talude não saturado sob chuva: SRM e GIM com e sem arc-length, sensibilidade e parâmetros \\
\multicolumn{3}{p{0.96\textwidth}}{\textbf{Formulação: campos aleatórios, FORM e amostragem por importância (docs/confiabilidade\_campos)}} \\
\path{docs/confiabilidade_campos/formulacao/formulacao.pdf} & 20 & o modelo, as duas bases de Karhunen--Loève (analítica e Nyström) e a mudança de base, o Monte Carlo, a AI, a variável de controle e o FORM com o gradiente modal (SRM e GIM); no apêndice, um exemplo de uma variável em Mathematica (exemplo\_1d.wl) \\
\path{docs/confiabilidade_campos/implementacao_form_ai/implementacao_form_ai.pdf} & 5 & como o FORM e a AI foram implementados no NeoPZ \\
\path{docs/confiabilidade_campos/ai_nystrom/ai_nystrom.pdf} & 5 & AI com a KL de Nyström, um roteiro testado em Mathematica \\
\path{docs/confiabilidade_campos/confiabilidade/confiabilidade.pdf} & 19 & FORM, Monte Carlo e AI em variáveis: a revisão de um código Mathematica e um talude infinito sob chuva (os números saem de Projects/Confiabilidade) \\
\path{docs/confiabilidade_campos/kl_form_placas/relatorio.pdf} & 26 & estudo do Ghanem \& Spanos (Seção 5.3): KL analítica, Nyström e Galerkin, Monte Carlo e FORM nas placas quadrada e curva \\
\path{docs/confiabilidade_campos/guia_estudo/guia_estudo.pdf} & 22 & guia de estudo que liga tudo (inclui o FORM com a base da KL, Eq. 14) \\
\multicolumn{3}{p{0.96\textwidth}}{\textbf{Resultados: o plano de trabalho (Projects/PlanoDeTrabalho)}} \\
\path{Projects/PlanoDeTrabalho/README.md} &  & o plano de trabalho (etapas 0 a 5, 2b e 2c), as regras dos documentos e este índice \\
\path{Projects/PlanoDeTrabalho/situacao/situacao.pdf} & --- & a situação: o plano, os resultados principais (com as seções do relatório das etapas), as hipóteses e os próximos passos (este) \\
\path{Projects/PlanoDeTrabalho/ETAPA0/relatorio_etapa0.pdf} & 9 & etapa 0: a AI da chuva no fermi e a prévia do Cho2010 (SRM antigo; superado pelo das etapas 1 a 3) \\
\path{Projects/PlanoDeTrabalho/ETAPAS1a3/relatorio_etapas.pdf} & 15 & todos os resultados das etapas 0 a 3 e a referência da chuva: SRM corrigido, AI contra o direto, teste linear, as duas bases, o gradiente modal (SRM e GIM), o viés, a variável de controle, a fase 4 parcial, o piloto e o custo \\
\path{Projects/PlanoDeTrabalho/arquivo/} &  & cópias datadas (AAAA-MM-DD\_nome.pdf), nunca editadas: 2026-11-17\_relatorio\_parcial.pdf \\
\multicolumn{3}{p{0.96\textwidth}}{\textbf{Resultados: as campanhas}} \\
\path{Projects/PlanoDeTrabalho/campanhas/relatorio_chuva.pdf} & 16 & Monte Carlo da chuva (Rahardjo); a Seção 7 é a AI \\
\path{Projects/PlanoDeTrabalho/campanhas/relatorio_rebaixamento.pdf} & 11 & rebaixamento do Ceron: Monte Carlo direto e AI (validação: 15,70\,\% contra 15,87\,\%) \\
\path{Projects/PlanoDeTrabalho/campanhas/relatorio_solos.pdf} & 7 & os três solos do Rahardjo sob chuva \\
\path{Projects/SlopeRandomFieldRefina/relatorio/relatorio.pdf} & 36 & refinamento por realização e Karhunen--Loève no talude do Ceron (2025); fica com o programa, porque os scripts dele e os do KLCilindroMika gravam nesta pasta \\
\multicolumn{3}{p{0.96\textwidth}}{\textbf{Fora do repositório}} \\
\path{~/Dropbox/UnsaturatedSoil/referencias/} &  & artigos de terceiros (Cho 2010, Ceron 2025, a tese do Santos, ...) \\
\path{~/projects/runs/} &  & os dados brutos das rodadas, uma pasta por campanha \\
\bottomrule
\end{longtable}
Como usar cada programa: o README.md da pasta dele, em Projects/<programa>/ (fica com o código).

\end{document}
